\section{Công nghệ Frontend}

Frontend là phần giao diện người dùng của hệ thống thương mại điện tử, chịu trách nhiệm hiển thị danh mục sản phẩm, xử lý tương tác mua hàng và tích hợp luồng camera phục vụ chức năng thử quần áo ảo. Đề tài lựa chọn Vue.js làm framework chính, kết hợp Nuxt.js để hỗ trợ Server-Side Rendering và TailwindCSS để xây dựng giao diện hiện đại, tương thích trên nhiều kích thước màn hình.

\subsection{Vue.js}
Vue.js là một framework JavaScript mã nguồn mở được phát triển bởi Evan You, dùng để xây dựng giao diện người dùng (UI) và ứng dụng đơn trang (SPA). Vue.js áp dụng kiến trúc component-based, cho phép tách giao diện thành các thành phần độc lập, có thể tái sử dụng và bảo trì dễ dàng. Cơ chế reactivity của Vue tự động đồng bộ dữ liệu (data) với giao diện thông qua hệ thống Virtual DOM, giúp tối ưu hiệu năng render. Trong đề tài, Vue.js được dùng để xây dựng các thành phần như danh sách sản phẩm, giỏ hàng, chi tiết đơn hàng và khung hiển thị thử đồ ảo \cite{vuejs_docs}.

\subsection{Nuxt.js}
Nuxt.js là framework cấp cao được xây dựng trên nền Vue.js, hỗ trợ đầy đủ các chế độ render gồm Server-Side Rendering (SSR), Static Site Generation (SSG) và Client-Side Rendering (CSR). SSR giúp trang web tải nhanh hơn và tối ưu cho công cụ tìm kiếm (SEO) -- yếu tố quan trọng đối với một website thương mại điện tử cần khả năng tiếp cận khách hàng qua Google. Nuxt.js cũng cung cấp cấu trúc thư mục quy ước (pages, components, layouts, middleware), tự động sinh route và tích hợp sẵn các tính năng như state management qua Pinia, plugins và composables, giúp rút ngắn thời gian phát triển \cite{nuxt_assets}.

\subsection{TailwindCSS}
TailwindCSS là một utility-first CSS framework cho phép xây dựng giao diện nhanh chóng thông qua các class tiện ích có sẵn, thay vì viết CSS tùy chỉnh từ đầu. Cách tiếp cận này giúp đảm bảo sự thống nhất về thiết kế trong toàn bộ ứng dụng và giảm đáng kể kích thước file CSS cuối cùng nhờ cơ chế tree-shaking, chỉ giữ lại các class thực sự được sử dụng \cite{tailwindcss_docs}.

\subsection{TypeScript}
TypeScript là ngôn ngữ lập trình mã nguồn mở do Microsoft phát 
triển, là tập cha (superset) của JavaScript với khả năng kiểm 
tra kiểu tĩnh (static typing) \cite{typescript_docs}. Trong đề 
tài, TypeScript được sử dụng xuyên suốt ở cả hai lớp: tầng 
giao diện người dùng (Nuxt.js, Vue.js) để đảm bảo tính nhất 
quán kiểu dữ liệu khi giao tiếp với API, và tầng AR (Lens 
Studio) để viết các script điều khiển thành phần thực tế tăng 
cường, giúp phát hiện lỗi sớm và nâng cao khả năng bảo trì 
mã nguồn.
